\documentclass[11pt,letterpaper]{article}
\usepackage[T1]{fontenc}
\usepackage[utf8]{inputenc}
\usepackage{mathpazo}
\usepackage[margin=1in,headheight=15pt]{geometry}
\usepackage{amsmath,amssymb,amsthm,mathtools}
\usepackage{microtype}
\usepackage{aliascnt}
\usepackage{booktabs,tabularx,array,longtable,multirow}
\usepackage{float}
\usepackage{enumitem}
\usepackage{xcolor}
\usepackage{fancyhdr}
\usepackage{graphicx}
\usepackage{tikz}
\usepackage{hyperref}
\usepackage[nameinlink,noabbrev]{cleveref}

\definecolor{ink}{RGB}{25,45,66}
\definecolor{accent}{RGB}{33,96,119}
\definecolor{warm}{RGB}{137,78,31}
\definecolor{lightgray}{gray}{0.94}
\definecolor{paleblue}{RGB}{237,246,249}
\hypersetup{
  colorlinks=true,
  linkcolor=accent,
  citecolor=accent,
  urlcolor=accent,
  pdftitle={Bounded-Width Power-Set Compression},
  pdfauthor={Research manuscript prepared for Vladimir Reshetnikov with ChatGPT},
  pdfsubject={Finite Boolean-lattice compression and choiceless infinite obstructions},
  pdfkeywords={Boolean lattice, width, Sperner theorem, Dilworth theorem, power set, axiom of choice, finite-to-one map, formalization}
}
\urlstyle{same}
\pagestyle{fancy}
\fancyhf{}
\fancyhead[L]{\small\textsc{Bounded-Width Power-Set Compression}}
\fancyhead[R]{\small Research manuscript}
\fancyfoot[C]{\thepage}
\renewcommand{\headrulewidth}{0.35pt}
\setlength{\parskip}{4pt}
\setlength{\parindent}{1.25em}
\setlength{\emergencystretch}{2em}
\setcounter{tocdepth}{2}
\setlist[enumerate]{itemsep=3pt,topsep=5pt,leftmargin=2.1em}
\setlist[itemize]{itemsep=3pt,topsep=5pt,leftmargin=1.8em}
\numberwithin{equation}{section}

\theoremstyle{plain}
\newtheorem{theorem}{Theorem}[section]
\newaliascnt{lemma}{theorem}
\newtheorem{lemma}[lemma]{Lemma}
\aliascntresetthe{lemma}
\newaliascnt{proposition}{theorem}
\newtheorem{proposition}[proposition]{Proposition}
\aliascntresetthe{proposition}
\newaliascnt{corollary}{theorem}
\newtheorem{corollary}[corollary]{Corollary}
\aliascntresetthe{corollary}
\theoremstyle{definition}
\newaliascnt{definition}{theorem}
\newtheorem{definition}[definition]{Definition}
\aliascntresetthe{definition}
\newaliascnt{problem}{theorem}
\newtheorem{problem}[problem]{Problem}
\aliascntresetthe{problem}
\crefname{problem}{Problem}{Problems}
\Crefname{problem}{Problem}{Problems}
\newaliascnt{conjecture}{theorem}
\newtheorem{conjecture}[conjecture]{Conjecture}
\aliascntresetthe{conjecture}
\theoremstyle{remark}
\newaliascnt{remark}{theorem}
\newtheorem{remark}[remark]{Remark}
\aliascntresetthe{remark}
\newaliascnt{example}{theorem}
\newtheorem{example}[example]{Example}
\aliascntresetthe{example}

\newcommand{\Pow}{\mathcal P}
\newcommand{\B}{\mathbb B}
\newcommand{\NN}{\mathbb N}
\newcommand{\ZZ}{\mathbb Z}
\newcommand{\pwidth}{\operatorname{wd}}
\newcommand{\im}{\operatorname{im}}
\newcommand{\ran}{\operatorname{ran}}
\newcommand{\dom}{\operatorname{dom}}
\newcommand{\Fin}{\operatorname{Fin}}
\newcommand{\SC}{\mathsf{SC}}
\newcommand{\ZF}{\mathsf{ZF}}
\newcommand{\ZFC}{\mathsf{ZFC}}
\newcommand{\AC}{\mathsf{AC}}
\newcommand{\bfto}{\preccurlyeq_{\mathrm{bfto}}}
\newcommand{\fto}{\preccurlyeq_{\mathrm{fto}}}
\newcommand{\repo}[1]{\texttt{\detokenize{#1}}}
\newcolumntype{L}[1]{>{\raggedright\arraybackslash}p{#1}}
\newcolumntype{Y}{>{\raggedright\arraybackslash}X}
\newcommand{\status}[1]{\par\noindent\colorbox{lightgray}{\parbox{\dimexpr\linewidth-2\fboxsep\relax}{#1}}\par}
\newcommand{\resultbox}[1]{\par\medskip\noindent\fcolorbox{accent}{paleblue}{\parbox{\dimexpr\linewidth-2\fboxsep-2\fboxrule\relax}{#1}}\par\medskip}

\begin{document}
\hypersetup{pageanchor=false}
\begin{titlepage}
\color{ink}
\noindent\textsc{Set theory \;\textbar\; Extremal combinatorics \;\textbar\; Formal mathematics}
\vspace{0.85in}
\begin{flushleft}
{\Huge\bfseries Bounded-Width\\[5pt] Power-Set Compression}\par
\vspace{0.22in}
{\Large Exact finite thresholds, choiceless infinite obstructions,\\[3pt]
and a defect extension of Glazer's chain-fiber theorem}\par
\vspace{0.28in}
{\large A finite/infinite research problem with proof certificates and a formalization route}\par
\end{flushleft}
\vspace{0.30in}
\noindent\rule{\linewidth}{0.7pt}
\vspace{0.18in}
\noindent Prepared for Vladimir Reshetnikov\\
Research manuscript prepared with ChatGPT\\
5 October 2026
\vfill
\status{\textbf{Research status.} The finite theorems are proved from Dilworth's and Sperner's theorems, with an explicit symmetric-chain certificate. The local-expansion, Dedekind-infinitude, and antichain-transfer theorems are proved in the manuscript; the transfer theorem imports Forster's published finite-to-one strengthening of Cantor's theorem. The full choiceless classification for width bound $r\geq2$ is posed as an open problem, not claimed solved. Historical priority for the new formulation and intermediate observations has not been established. This manuscript is AI-assisted, unrefereed, and not formally verified. Elliot Glazer is a cited researcher, not an author or endorser.}
\vspace{7pt}
{\footnotesize\noindent\textbf{Repository snapshot.} The proposed integration targets the public \emph{ProveIt} repository at commit \texttt{b71545625}, inspected on 5 October 2026. The archive accompanying this manuscript contains the \LaTeX{} source, compiled PDF, and an exact standard-library Python verifier.\par}
\end{titlepage}
\hypersetup{pageanchor=true}

\pagenumbering{roman}
\section*{Abstract}
\addcontentsline{toc}{section}{Abstract}
For a poset $P$ and a positive integer $r$, call a map $c:P\to Y$ an
$r$-width compression when every fiber has width at most $r$.  The motivating
infinite question asks for such a map
\[
        f:\Pow(X)\longrightarrow X,
\]
where $\Pow(X)$ is ordered by inclusion.  At $r=1$ this is the chain-fiber
problem recently settled in $\ZF$ by Elliot Glazer: such a map exists exactly
when $1\leq |X|\leq3$.

We formulate the bounded-width extension and solve its finite form exactly.
For every finite poset $P$ of width $w$ and every list of fiber budgets
$r_1,\dots,r_q$, a coloring with $i$th fiber of width at most $r_i$ exists if
and only if $\sum_i r_i\geq w$.  Hence, for the Boolean lattice
$\B_n=\Pow([n])$,
\[
 q_r(n)=\min\{|Y|:\B_n\to Y\text{ has fibers of width }\leq r\}
       =\left\lceil\frac{\binom n{\lfloor n/2\rfloor}}{r}\right\rceil.
\]
A self-indexed compression $\Pow([n])\to[n]$ therefore exists exactly when
$\binom n{\lfloor n/2\rfloor}\leq rn$.  Its least admissible defect is
\[
 \delta(n)=\left\lceil\frac1n\binom n{\lfloor n/2\rfloor}\right\rceil
 \sim \sqrt{\frac2\pi}\,\frac{2^n}{n^{3/2}},
\]
and the largest feasible $n$ for a fixed $r$ is
\[
 N(r)=\log_2 r+\frac32\log_2\log_2 r
       +\frac12\log_2\!\frac\pi2+O(1).
\]
The upper bound is constructive: group the chains in a symmetric-chain
decomposition, at most $r$ chains per label.  The middle layer is the matching
lower-bound certificate.

For the infinite form we prove three choice-free results.  First, every finite
$S\subseteq X$ obeys the exact expansion inequality
\[
 |f[\Pow(S)]|\geq q_r(|S|).
\]
Consequently every hypothetical infinite example canonically generates a
countably infinite subset by iterating $S\mapsto S\cup f[\Pow(S)]$; in
particular no Dedekind-finite or amorphous set can be a counterexample.
Second, an antichain-transfer theorem converts any antichain code
$\Pow(U)\to\Pow(X)$ together with a finite-to-one map $X\to U$ into a
finite-to-one map $\Pow(U)\to U$, contradicting Forster when $U$ is infinite.
Third, if $2\times X$ injects into $X$, the complementary-transversal code
supplies such an antichain.  Thus no well-orderable infinite $X$ works, and
under Choice the finite classification is complete.  The unresolved $\ZF$
frontier for $r\geq2$ is sharply localized: an infinite counterexample would
have to be Dedekind-infinite but non-duplicable, with no suitable
finite-to-one core.

\vspace{0.12in}
\noindent\textbf{Keywords.} Boolean lattice; poset width; Dilworth theorem; Sperner theorem; symmetric-chain decomposition; power set; axiom of choice; finite-to-one map; Dedekind-infinite set; formal proof certificate.

\clearpage
\begingroup
\fontsize{10}{11}\selectfont
\setlength{\parskip}{0pt}
\tableofcontents
\endgroup
\clearpage
\pagenumbering{arabic}

\section{The problem and why the finite/infinite bridge matters}\label{sec:intro}

\subsection{The chain-fiber starting point}
In December 2025 Alexander Pruss asked whether $\ZF$ alone rules out, for an
infinite set $X$, a map $f:\Pow(X)\to X$ that separates every incomparable
pair: $f(A)\neq f(B)$ whenever neither $A\subseteq B$ nor $B\subseteq A$.
Equivalently, each fiber $f^{-1}(x)$ must be a chain under inclusion.  Elliot
Glazer answered with a sharp theorem: in $\ZF$, the map exists exactly for
sets of cardinality $1$, $2$, or $3$ \cite{PrussGlazer2025}.  Emil
Je\v{r}\'abek supplied a complementary cardinal-arithmetic argument on the
same thread.

That result has an unusually clean finite/infinite profile.  The finite cutoff
is not merely an artifact of cardinality: it is the width of a Boolean
lattice.  The infinite proof, by contrast, must survive the absence of a
well-order and of the finite choices that would normally split a bounded-width
poset into chains.  This suggests varying the amount of incomparability
allowed inside one label.

\subsection{The bounded-width extension}
Fix a positive integer $r$.  We permit each label to identify several
mutually incompatible subsets, but cap the size of an antichain inside one
fiber by $r$.

\begin{problem}[bounded-width power-set compression]\label{prob:main}
Classify the pairs $(X,r)$ for which there is a function
\[
 f:\Pow(X)\longrightarrow X
\]
such that every fiber $f^{-1}(x)$ has inclusion-width at most $r$.
Determine the exact finite threshold, the strongest theorem provable in
$\ZF$, and the additional strength supplied by choice principles.
\end{problem}

The parameter $r$ is a \emph{contradiction budget}: a label may be attached to
at most $r$ pairwise incomparable descriptions.  The case $r=1$ is Glazer's
theorem.  For $r>1$ the problem retains a fully explicit finite extremal
answer but opens a genuinely choiceless infinite frontier.

\subsection{Why this should interest Glazer}
The problem directly extends a theorem proved by Glazer, while its infinite
part lies at the intersection of his documented interests in choiceless set
theory and hard, independently checkable mathematical problems.  His work
with Bokai Yao separates principles that coincide under Choice but diverge in
$\ZF$ with urelements \cite{GlazerYao2026}; the present problem similarly has
a complete $\ZFC$ classification and a narrower unresolved $\ZF$ boundary.
His work on FrontierMath emphasizes original expert-level problems with
objective verification surfaces \cite{FrontierMath}; here the finite theorem
has compact upper and lower certificates, while the infinite theorem has a
precise imported dependency and a transparent open remainder.

\subsection{Relation to ProveIt}
The current ProveIt repository describes itself as machine-checked mathematics
in Lean and Rocq/Coq, with executable certificates separated from untrusted
exploratory computations \cite{ProveIt}.  This problem naturally spans three
of its active surfaces:
\begin{enumerate}
\item finite posets and Boolean-lattice certificates in combinatorics;
\item first-order $\ZF$ and cardinal arithmetic in set theory;
\item a small proof-producing checker whose trusted boundary is simpler than
      a general-purpose optimizer.
\end{enumerate}
The finite core is suitable for immediate formalization.  The infinite side
suggests a staged project: formalize the local expansion and antichain code
first, then isolate Forster's theorem as the only substantial imported
set-theoretic component.

\subsection{Contribution boundary}
The exact finite formula is a short synthesis of classical results of
Dilworth and Sperner, enhanced here by a nonuniform budget theorem, an
explicit certificate format, threshold inversion, and its connection to the
choiceless problem.  The local-expansion, definable Dedekind-infinitude, and
general antichain-transfer statements are proved below.  A targeted search of
the inspected literature did not locate this precise bounded-width
formulation, but absence from a search is not a priority proof.  No claim is
made that the unresolved $r\geq2$ question is known to be open in every
source not inspected.

\resultbox{\textbf{Main outcome.}  The finite problem and the entire $\ZFC$
classification are solved exactly.  In $\ZF$, every possible infinite
counterexample for $r\geq2$ is forced into a narrow cardinal-arithmetic gap:
it is Dedekind-infinite, it is not binary-duplicable, and it has no
finite-to-one core carrying two disjoint copies of itself.}

\section{Definitions and the classification targets}\label{sec:definitions}

\begin{definition}[width]\label{def:width}
For a poset $P$, its width $\pwidth(P)$ is the supremum of the cardinalities of
its antichains.  When $r\in\NN_{>0}$, the assertion $\pwidth(P)\leq r$ means that
$P$ has no antichain of size $r+1$.
\end{definition}

\begin{definition}[width compression]\label{def:compression}
Let $P$ be a poset, let $Y$ be a set, and let $r\geq1$.  An
\emph{$r$-width compression} is a function $c:P\to Y$ such that
\[
       \pwidth(c^{-1}(y))\leq r \qquad(y\in Y).
\]
For a finite poset define
\[
 \chi_r(P)=\min\{|Y|: \text{an $r$-width compression }P\to Y\text{ exists}\}.
\]
\end{definition}

\begin{definition}[Boolean and self-indexed parameters]\label{def:booleanparams}
Write $\B_n=(\Pow([n]),\subseteq)$ and
\[
 W_n=\pwidth(\B_n)=\binom n{\lfloor n/2\rfloor}.
\]
Set
\[
 q_r(n)=\chi_r(\B_n),\qquad
 \delta(n)=\min\{r:\Pow([n])\to[n]\text{ is an $r$-width compression}\},
\]
for $n\geq1$, and
\[
 N(r)=\max\{n\geq1:\delta(n)\leq r\}.
\]
\end{definition}

For arbitrary sets it is useful to name the existence statement.

\begin{definition}[self-compression statement]\label{def:SC}
For $r\geq1$, let $\SC_r(X)$ denote:
\[
 \exists f:\Pow(X)\to X\quad
 \forall x\in X\quad \pwidth(f^{-1}(x))\leq r.
\]
\end{definition}

The strongest natural classification suggested by the finite answer is the
following.

\begin{conjecture}[bounded-width Glazer classification]\label{conj:BWG}
For every fixed $r\geq1$, $\ZF$ proves that $\SC_r(X)$ holds if and only if
$X$ is finite, nonempty, and
\[
 \binom{|X|}{\lfloor |X|/2\rfloor}\leq r|X|.
\]
\end{conjecture}

For $r=1$ this is Glazer's theorem.  We prove the finite direction in $\ZF$,
the full statement in $\ZFC$, and several choice-free restrictions on a
hypothetical counterexample for $r\geq2$.

\section{A universal finite theorem with nonuniform budgets}\label{sec:finiteposets}

The first result is not specific to Boolean lattices.  It identifies width as
the complete resource controlling finite compression.

\begin{theorem}[budgeted finite-poset theorem]\label{thm:budget}
Let $P$ be a finite poset of width $w$, and let $r_1,\dots,r_q$ be
nonnegative integers.  There is a map $c:P\to[q]$ satisfying
\[
       \pwidth(c^{-1}(i))\leq r_i\qquad(1\leq i\leq q)
\]
if and only if
\[
       \sum_{i=1}^q r_i\geq w.
\]
\end{theorem}

\begin{proof}
Let $A$ be an antichain of size $w$.  Each fiber contains at most $r_i$
points of $A$, so
\[
 w=|A|=\sum_{i=1}^q |A\cap c^{-1}(i)|\leq\sum_{i=1}^q r_i.
\]
This proves necessity.

For sufficiency, Dilworth's decomposition theorem partitions $P$ into exactly
$w$ chains \cite{Dilworth1950}.  Distribute these chains among bins
$1,\dots,q$, placing at most $r_i$ chains in bin $i$; the capacity inequality
makes this possible.  Color every point by the bin containing its chain.  An
antichain meets each chain at most once, so the union of at most $r_i$ chains
has width at most $r_i$.
\end{proof}

\begin{corollary}[uniform compression number]\label{cor:uniform}
For every finite poset $P$ and every $r\geq1$,
\[
              \chi_r(P)=\left\lceil\frac{\pwidth(P)}r\right\rceil.
\]
\end{corollary}

\begin{proof}
Apply \cref{thm:budget} with all capacities equal to $r$.
\end{proof}

\begin{remark}[partition versus cover]\label{rem:cover}
The same criterion holds if one begins with a cover by width-bounded
subposets.  Choose one covering set for each point and delete the point from
all other sets; taking subposets cannot increase width.  Thus the function
and cover formulations are equivalent in the finite theorem.
\end{remark}

\begin{remark}[why the proof does not settle the infinite problem]
Two finite ingredients become delicate without Choice.  A finite-width
infinite poset need not come equipped with a selectable decomposition into
finitely many chains, and even if each individual fiber can be decomposed,
choosing coherent decompositions for an $X$-indexed family is itself a choice
problem.  More importantly, replacing $X$ by $X\times r$ can change
cardinality in $\ZF$.
\end{remark}

\section{Boolean lattices: exact formula and explicit certificates}\label{sec:boolean}

Sperner's theorem gives $\pwidth(\B_n)=W_n$ \cite{Sperner1928}.  Combining it with
\cref{cor:uniform} yields the exact finite answer.

\begin{theorem}[exact Boolean-lattice compression]\label{thm:booleanexact}
For all $n\geq0$ and $r\geq1$,
\[
       q_r(n)=\left\lceil\frac{W_n}{r}\right\rceil
             =\left\lceil\frac1r\binom n{\lfloor n/2\rfloor}\right\rceil.
\]
\end{theorem}

The proof has a particularly transparent pair of certificates.

\subsection{Lower certificate: the middle layer}
The family
\[
       L_n=\{A\subseteq[n]:|A|=\lfloor n/2\rfloor\}
\]
is an antichain of size $W_n$.  A color can contain at most $r$ members of
$L_n$, so at least $\lceil W_n/r\rceil$ colors are required.  The certificate
is simply the list of all middle-layer subsets.

\subsection{Upper certificate: grouped symmetric chains}
The Boolean lattice has a symmetric-chain decomposition: a partition into
saturated chains whose bottom and top ranks sum to $n$ \cite{deBruijn1951,
GreeneKleitman1976}.  Every such decomposition has exactly $W_n$ chains,
because each chain meets rank $\lfloor n/2\rfloor$ exactly once.  Group at
most $r$ chains under each label.  Every fiber is then a union of at most
$r$ chains and consequently has width at most $r$.

A recursive decomposition can be constructed without optimization.  Begin
with the sole chain $\varnothing$ in $\B_0$.  Suppose
\[
 C:A_k\subset A_{k+1}\subset\cdots\subset A_{n-1-k}
\]
is a symmetric chain in $\B_{n-1}$.  Replace it in $\B_n$ by
\begin{align*}
 C^+&:A_k\subset\cdots\subset A_{n-1-k}
       \subset A_{n-1-k}\cup\{n\},\\
 C^-&:A_k\cup\{n\}\subset\cdots\subset A_{n-2-k}\cup\{n\},
\end{align*}
where the second chain is omitted when empty.  The two new chains are
symmetric, disjoint, and together contain both copies of every old member.
Induction gives a full decomposition.

\begin{example}[$\B_4$]\label{ex:B4}
The recursive construction gives six chains:
\begin{align*}
C_1&:\varnothing<1<12<123<1234, &
C_2&:4<14<124,\\
C_3&:3<13<134,                 &
C_4&:34,\\
C_5&:2<23<234,                 &
C_6&:24.
\end{align*}
Here a string denotes the corresponding subset.  For $r=2$, the groups
$(C_1,C_2)$, $(C_3,C_4)$, and $(C_5,C_6)$ give a three-label compression.
The six two-subsets form the matching lower certificate, so three is optimal.
\end{example}

\subsection{A proof-certificate interface}
An upper certificate consists of a list of chains and a group number for each
chain.  A checker verifies:
\begin{enumerate}
\item every element of $\B_n$ occurs exactly once;
\item consecutive members of each chain differ by adding one element;
\item each group contains at most $r$ chains.
\end{enumerate}
The checker never computes a width.  The lower certificate is the middle
layer and the integer inequality $|L_n|>r(q-1)$.  Thus optimality can be
kernel-checked by elementary finite-set operations.

\section{Self-indexed finite compression}\label{sec:selfindexed}

A map $\Pow([n])\to[n]$ may leave labels unused.  Therefore it exists exactly
when the minimum number of labels is at most $n$.

\begin{theorem}[complete finite classification]\label{thm:finiteclass}
Let $n\geq1$ and $r\geq1$.  The following are equivalent:
\begin{enumerate}
\item there is an $r$-width compression $\Pow([n])\to[n]$;
\item $q_r(n)\leq n$;
\item $W_n\leq rn$;
\item $r\geq\delta(n)$, where
\[
 \delta(n)=\left\lceil\frac{W_n}{n}\right\rceil.
\]
\end{enumerate}
\end{theorem}

\begin{proof}
Only the equivalence between (1) and (2) needs comment.  A map to $[n]$ uses
at most $n$ nonempty fibers, so (1) implies $q_r(n)\leq n$.  Conversely, an
optimal compression with $q_r(n)\leq n$ labels can inject its label set into
$[n]$; the remaining labels have empty fibers.
\end{proof}

\begin{corollary}[the first defects]\label{cor:firstdefects}
The least width budgets are
\[
\delta(1),\delta(2),\dots=
1,1,1,2,2,4,5,9,14,26,42,77,132,246,429,805,\dots.
\]
In particular:
\begin{enumerate}
\item $r=1$ permits exactly $n=1,2,3$;
\item $r=2$ permits exactly $1\leq n\leq5$;
\item $r=3$ still permits only $1\leq n\leq5$;
\item $r=4$ permits exactly $1\leq n\leq6$.
\end{enumerate}
\end{corollary}

The finite $r=1$ line coincides with Glazer's full $\ZF$ theorem, which is a
useful consistency check on the formulation.

\subsection{Monotonicity}
The threshold is genuinely one-sided; no later value of $n$ becomes feasible
after failure.

\begin{proposition}[monotonic defect]\label{prop:monotonic}
The sequence $W_n/n$ is nondecreasing for $n\geq1$.  Hence $\delta(n)$ is
nondecreasing and
\[
 N(r)=\max\{n:W_n\leq rn\}
\]
is well-defined.
\end{proposition}

\begin{proof}
Put $a_n=W_n/n$.  For $m\geq1$,
\[
 \frac{a_{2m+1}}{a_{2m}}
 =\frac{2m}{m+1}\geq1,
\qquad
 \frac{a_{2m+2}}{a_{2m+1}}
 =\frac{2m+1}{m+1}>1.
\]
Also $a_1=a_2=1$.  Taking ceilings preserves monotonicity.  Finally,
$W_n\geq2^n/(n+1)$ because the largest binomial coefficient is at least the
average of the $n+1$ coefficients.  Thus $W_n/n\to\infty$, so the defining
set for $N(r)$ is finite.
\end{proof}

\begin{remark}[plateaus]
Ceiling creates a few initial plateaus: $\delta(1)=\delta(2)=\delta(3)=1$ and
$\delta(4)=\delta(5)=2$.  Thereafter the values grow quickly.  The real ratio
$W_n/n$ itself is strictly increasing except for the first two parity steps
where the displayed ratios equal one.
\end{remark}

\section{Asymptotic defect and inversion}\label{sec:asymptotics}

The exact formula makes the growth transparent.  Stirling's expansion,
uniform across the two parities, gives
\begin{equation}\label{eq:centralasymp}
 W_n=2^n\sqrt{\frac{2}{\pi n}}\left(1+O(n^{-1})\right).
\end{equation}

\begin{theorem}[defect asymptotic]\label{thm:defectasymp}
As $n\to\infty$,
\[
 \delta(n)=\sqrt{\frac2\pi}\,\frac{2^n}{n^{3/2}}
              \left(1+O(n^{-1})\right).
\]
\end{theorem}

\begin{proof}
Divide \cref{eq:centralasymp} by $n$.  The ratio tends to infinity, so taking
the ceiling changes it by a relative $o(1)$ amount, absorbed by the stated
error.
\end{proof}

Thus a fixed multiplicative allowance in fiber width buys only an additive
amount in the universe size.  More precisely:

\begin{theorem}[threshold inversion]\label{thm:thresholdasymp}
Let $L=\log_2 r$.  As $r\to\infty$,
\[
 N(r)=L+\frac32\log_2 L+\frac12\log_2\!\frac\pi2+O(1).
\]
\end{theorem}

\begin{proof}
At the smooth transition $W_n/n\asymp r$, \cref{eq:centralasymp} gives
\[
 n-\frac32\log_2 n+\frac12\log_2\frac2\pi
      =L+O(n^{-1}).
\]
The equation first implies $n=L+O(\log L)$.  Substituting this estimate inside
$\log_2 n$ yields
\[
 n=L+\frac32\log_2 L+\frac12\log_2\frac\pi2
       +O\!\left(\frac{\log L}{L}\right)
\]
for the real transition.  The integer maximum $N(r)$ differs from that
transition by at most a bounded rounding and parity term, which gives the
stated $O(1)$ remainder.
\end{proof}

\begin{remark}[interpretation]
The factor $n^{-3/2}$ has two sources: $n^{-1/2}$ from concentration of the
binomial layers and another $n^{-1}$ because the available label set has only
$n$ elements.  Inverting $2^n/n^{3/2}$ creates the secondary
$\tfrac32\log_2\log_2 r$ term.
\end{remark}

\section{The infinite problem in \texorpdfstring{$\ZF$}{ZF}}\label{sec:infinite}

\subsection{The solved chain case}
The following theorem is the exact $r=1$ endpoint motivating the present
extension.

\begin{theorem}[Glazer, 2025]\label{thm:glazer}
$\ZF$ proves, for every set $X$, that the following are equivalent:
\begin{enumerate}
\item $1\leq|X|\leq3$;
\item there is $f:\Pow(X)\to X$ whose fibers are inclusion-chains;
\item $\Pow(X)$ is the union of an $X$-indexed family of inclusion-chains.
\end{enumerate}
\end{theorem}

The published form here is a MathOverflow answer rather than a journal paper
\cite{PrussGlazer2025}.  The proof is choice-free and uses a diagonal
construction after extracting enough coding power from the two-element layer.
It is stronger than merely observing that the finite width $W_n$ exceeds $n$
for $n\geq4$.

\subsection{Why \texorpdfstring{$r\geq2$}{r >= 2} is structurally different}
For one chain, each rank contains at most one member, and the two-element
layer supplies a single-valued coding object.  A width-$r$ fiber can contain
$r$ members of the same rank.  Turning those finite sets into $r$ separately
named chains requires choosing an ordering or coloring inside many finite
fibers.  Such finite choices cannot be silently made in $\ZF$.

This is exactly the setting in which boundedly finite-to-one maps become
subtle.  Hu and Shen emphasize that, without Choice, a boundedly finite-to-one
map need not yield an injection; they prove several generalized Cantor
nonexistence theorems while recording models where familiar implications fail
\cite{HuShen2025}.  Shen's work on finitary partitions likewise shows that a
natural Cantor analogue can consistently fail when the power-set functor is
replaced by another large combinatorial construction \cite{Shen2026}.

The next two sections isolate what can nevertheless be proved directly from
the order structure of $\Pow(X)$.

\section{Finite expansion inside every infinite candidate}\label{sec:expansion}

The finite theorem applies to every finite subcube, even though its labels may
land anywhere in $X$.

\begin{theorem}[exact local expansion]\label{thm:localexpansion}
Assume $\SC_r(X)$, witnessed by $f:\Pow(X)\to X$.  For every finite
$S\subseteq X$ with $|S|=n$,
\[
       |f[\Pow(S)]|\geq q_r(n)
       =\left\lceil\frac{W_n}{r}\right\rceil.
\]
Consequently,
\[
 |f[\Pow(S)]\setminus S|\geq
       \max\left(0,\left\lceil\frac{W_n}{r}\right\rceil-n\right).
\]
\end{theorem}

\begin{proof}
Restrict $f$ to the finite Boolean lattice $\Pow(S)$.  For every label $x$,
the restricted fiber $f^{-1}(x)\cap\Pow(S)$ is a subposet of the global
fiber, hence has width at most $r$.  By \cref{thm:booleanexact}, the restricted
map uses at least $q_r(n)$ distinct labels.  At most $n$ of those labels lie
in $S$, giving the second inequality.
\end{proof}

\begin{corollary}[no large finite closed set]\label{cor:nofiniteclosed}
If $n>N(r)$, no $n$-element set $S\subseteq X$ can satisfy
$f[\Pow(S)]\subseteq S$.
\end{corollary}

\begin{proof}
For $n>N(r)$, $q_r(n)>n$ by definition.  The image cannot both have more than
$n$ elements and be contained in $S$.
\end{proof}

Define the finite-hull operator
\[
             H_f(S)=S\cup f[\Pow(S)].
\]
Then \cref{thm:localexpansion} gives the exact recurrence
\begin{equation}\label{eq:hullgrowth}
       |H_f(S)|\geq q_r(|S|)
       \qquad\text{whenever }|S|>N(r).
\end{equation}
Since $W_n\geq2^n/(n+1)$, this eventually dominates
$2^n/(r(n+1))$; iterating the hull therefore grows faster than an ordinary
exponential recurrence.

\begin{theorem}[definable Dedekind-infinitude]\label{thm:dedekind}
Fix $r\geq1$.  In $\ZF$, if $X$ is infinite and $\SC_r(X)$ holds, then $X$
is Dedekind-infinite.  More precisely, from a witnessing map $f$ and any
ordered $(N(r)+1)$-tuple of distinct elements of $X$, one can define an
injection $\NN\to X$.
\end{theorem}

\begin{proof}
Let $S_0\subseteq X$ have $N(r)+1$ elements, equipped with the order supplied
by the chosen tuple.  Recursively put
\[
       S_{k+1}=H_f(S_k)=S_k\cup f[\Pow(S_k)].
\]
Each $S_k$ is finite.  Since $|S_k|>N(r)$, \cref{eq:hullgrowth} gives
$|S_{k+1}|>|S_k|$.

It remains to choose a new point at every stage without invoking countable
choice.  We define an order on $S_{k+1}$ canonically from the order on $S_k$.
The latter induces a lexicographic order on the finite set $\Pow(S_k)$.  For
$y\in f[\Pow(S_k)]$, let $\mu_k(y)$ be the least subset $A\subseteq S_k$ with
$f(A)=y$.  The map $\mu_k$ is injective, so it transports the lexicographic
order to a linear order on $f[\Pow(S_k)]$.  Order $S_{k+1}$ by listing the old
points in their old order, followed by the new points in this transported
order.

Let $x_k$ be the first point of $S_{k+1}\setminus S_k$ in that order.  The
construction is recursive and functional, so Replacement produces the
sequence $k\mapsto x_k$.  The points are pairwise distinct because
$x_k\in S_{k+1}\setminus S_k$ and the $S_k$ are increasing.  Hence this
sequence is an injection from $\NN$ into $X$.
\end{proof}

\begin{corollary}[Dedekind-finite obstruction]\label{cor:dedekindfinite}
No infinite Dedekind-finite set, and in particular no amorphous set, satisfies
$\SC_r(X)$ for any finite $r$.
\end{corollary}

\begin{remark}[more than existential infinitude]
The proof does not appeal to the principle that a countable union of nonempty
finite sets has a choice function.  It manufactures a canonical local order
from least preimages at each stage.  Relative to $f$ and the finite seed, the
countable subset is definable.  This feature is useful for formalization and
may be stronger than the bare Dedekind-infinitude conclusion in symmetric
models.
\end{remark}

\section{Antichain transfer and the duplicability obstruction}\label{sec:transfer}

The second infinite mechanism extracts a finite-to-one map from a large,
explicit antichain.

\begin{definition}[antichain code]\label{def:antichaincode}
An \emph{antichain code from $U$ into $X$} is a map
\[
        \Theta:\Pow(U)\longrightarrow\Pow(X)
\]
whose range is an antichain under inclusion.
\end{definition}

\begin{lemma}[width becomes bounded multiplicity]\label{lem:widthmultiplicity}
If $f:\Pow(X)\to X$ has fibers of width at most $r$ and $\Theta$ is an
antichain code from $U$ into $X$, then $f\circ\Theta:\Pow(U)\to X$ is at most
$r$-to-one.
\end{lemma}

\begin{proof}
If $r+1$ distinct subsets of $U$ had the same image, their $\Theta$-images
would be $r+1$ pairwise incomparable members of one fiber of $f$.
\end{proof}

We use Forster's theorem in its published $\ZF$ form.

\begin{theorem}[Forster]\label{thm:forster}
If there is a finite-to-one map $\Pow(U)\to U$, then $U$ is finite.
\end{theorem}

\noindent Forster's result strictly strengthens Cantor's theorem in the
absence of Choice \cite{Forster2003}.

\begin{theorem}[antichain-transfer obstruction]\label{thm:transfer}
Let $U$ be infinite.  Suppose there are
\begin{enumerate}
\item an antichain code $\Theta:\Pow(U)\to\Pow(X)$, and
\item a finite-to-one map $p:X\to U$.
\end{enumerate}
Then $\SC_r(X)$ fails for every finite $r$.
\end{theorem}

\begin{proof}
If $f$ witnessed $\SC_r(X)$, then $h=f\circ\Theta$ would be at most
$r$-to-one by \cref{lem:widthmultiplicity}.  For a point $u\in U$, the fiber
of $p\circ h$ is the union, over the finite set $p^{-1}(u)$, of fibers of
$h$ of size at most $r$.  Hence $p\circ h:\Pow(U)\to U$ is finite-to-one,
contradicting \cref{thm:forster}.
\end{proof}

The required antichain code has a simple two-channel construction.

\begin{lemma}[complementary-transversal code]\label{lem:transversal}
If there is an injection $e:U\times2\to X$, then there is an antichain code
$\Theta:\Pow(U)\to\Pow(X)$.
\end{lemma}

\begin{proof}
Write $e_i(u)=e(u,i)$.  The two ranges are disjoint.  Define
\[
       \Theta(A)=e_0[A]\cup e_1[U\setminus A].
\]
If $A\neq B$, choose $u$ in their symmetric difference.  If
$u\in A\setminus B$, then $e_0(u)\in\Theta(A)\setminus\Theta(B)$ while
$e_1(u)\in\Theta(B)\setminus\Theta(A)$.  Thus neither coded set contains the
other.  The other case is symmetric.
\end{proof}

\begin{theorem}[two-channel core criterion]\label{thm:core}
Let $U$ be infinite.  If
\[
             U\times2\preccurlyeq X
       \quad\text{and}\quad
             X\fto U,
\]
then $\SC_r(X)$ fails for every finite $r$.
\end{theorem}

\begin{proof}
The first map supplies the antichain code by \cref{lem:transversal}; the
second map and \cref{thm:transfer} finish the proof.
\end{proof}

\begin{corollary}[binary duplicability obstruction]\label{cor:duplicate}
If $X$ is infinite and $X\times2\preccurlyeq X$, then $\SC_r(X)$ fails for
every finite $r$.
\end{corollary}

\begin{proof}
Use \cref{thm:core} with $U=X$ and $p=\operatorname{id}_X$.
\end{proof}

\begin{corollary}[well-orderable and choice cases]\label{cor:choice}
No infinite well-orderable set satisfies $\SC_r(X)$.  Consequently $\ZFC$
proves that $\SC_r(X)$ holds exactly for the finite nonempty sets satisfying
$W_{|X|}\leq r|X|$.
\end{corollary}

\begin{proof}
Every infinite well-orderable set is equipotent with an infinite initial
ordinal $\kappa$, and $2\times\kappa$ injects into $\kappa$.  Apply
\cref{cor:duplicate}.  Under $\AC$, every set is well-orderable.  The finite
cases are exactly \cref{thm:finiteclass}.
\end{proof}

\begin{remark}[why Forster is the right dependency]
Cantor's theorem alone excludes an injection $\Pow(U)\to U$, but
$f\circ\Theta$ is only $r$-to-one.  Turning that map into an injection would
silently use a finite choice operation on its fibers.  Forster's theorem is
precisely the choice-free statement needed here.
\end{remark}

\section{The remaining choiceless frontier}\label{sec:frontier}

Combining the preceding sections yields a sharp profile of any counterexample
to \cref{conj:BWG}.

\begin{theorem}[counterexample profile]\label{thm:profile}
Fix $r\geq2$.  If an infinite set $X$ satisfies $\SC_r(X)$, then:
\begin{enumerate}
\item $X$ is Dedekind-infinite, and a countable subset is definable from the
      witness and finitely many parameters;
\item $X\times2\not\preccurlyeq X$;
\item more generally, there is no infinite $U$ with
      $U\times2\preccurlyeq X$ and $X\fto U$;
\item every finite $S\subseteq X$ of size $n$ emits at least
      $q_r(n)-n$ labels outside itself whenever this number is positive;
\item the ambient model violates Choice.
\end{enumerate}
\end{theorem}

\begin{proof}
Items (1)--(4) are \cref{thm:dedekind,thm:core,thm:localexpansion}.  Item (5)
follows from \cref{cor:choice}.
\end{proof}

This profile rules out the most familiar choiceless candidates.  An amorphous
set is impossible by (1); an infinite cardinal is impossible by (2); and a
set that is finite-to-one reducible to an infinite two-channel core is
impossible by (3).  The only remaining region consists of Dedekind-infinite
sets with highly nonclassical finite-multiple arithmetic.

\subsection{Why the obvious reductions stop here}
One might try to decompose each width-$r$ fiber into $r$ chains and apply the
$r=1$ theorem.  There are two obstructions.
\begin{enumerate}
\item Infinite Dilworth decompositions can require compactness or choice.
\item Even if each fiber separately admits a decomposition, choosing one for
      every $x\in X$ may be a family-choice principle; afterward one still
      has an $X\times r$-indexed family, and $X\times r$ need not reduce to
      $X$ in $\ZF$.
\end{enumerate}
Likewise, restricting $f$ to one rank produces an $r$-to-one map from a
finite-subset functor into $X$, but boundedly finite-to-one maps do not in
general canonically split into injections without Choice \cite{HuShen2025}.

\subsection{Two credible routes forward}
A proof of \cref{conj:BWG} could seek a \emph{canonical} chain splitting
special to families of subsets, avoiding arbitrary finite choices.  A
counterexample would likely require a permutation or symmetric model with a
Dedekind-infinite, non-duplicable set $X$ and an equivariant map on its full
power set.  The local expansion theorem imposes a severe support condition:
large finite subcubes must acquire exponentially many distinct output atoms.
Either route is concrete enough for a focused research program.

\section{Algorithms, exact checks, and certificate complexity}\label{sec:checks}

The accompanying script \repo{verification.py} uses only Python's standard
library.  Its role is diagnostic, not foundational: the proofs above are
symbolic.

\subsection{What was checked}
The executed verifier performs the following exact checks.
\begin{enumerate}
\item It constructs the recursive symmetric-chain decomposition for
      $0\leq n\leq14$ and verifies coverage, disjointness, saturation,
      symmetry, and the chain count $W_n$.
\item For $1\leq n\leq10$ and $1\leq r\leq5$, it groups chains and verifies
      that the number of groups is $\lceil W_n/r\rceil$.
\item It checks exact rational monotonicity of $W_n/n$ and monotonicity of
      $\delta(n)$ for $1\leq n\leq300$.
\item It generates the defect and threshold tables in this manuscript and a
      numerical diagnostic for \cref{thm:thresholdasymp}.
\end{enumerate}
All assertions passed.  The full textual output is included as
\repo{verification_report.txt}.

\subsection{Complexity}
An explicit decomposition necessarily names all $2^n$ subsets, so its output
size is exponential.  The verifier runs in time polynomial in that output
size.  The lower certificate is much smaller only in representation, not in
cardinality: it contains $W_n\asymp2^n/\sqrt n$ middle-layer subsets.

For formal work, one can avoid materializing the entire certificate by
proving the recursive SCD algorithm correct once and evaluating only the
closed formula for the number of chains.  This separates an executable
instance certificate from a reusable theorem.

\subsection{Pseudocode}
\begin{verbatim}
SCD(0) := [[empty set]]
for new = 1,...,n:
    output := []
    for chain C = [A_k,...,A_t] in SCD(new-1):
        append [A_k,...,A_t,A_t union {new}] to output
        if C has at least two members:
            append [A_k union {new},...,A_{t-1} union {new}] to output
    chains := output
return chains
\end{verbatim}
The two output chains partition the two copies of the input chain in
$\B_{n-1}\times\{0,1\}$, which makes the inductive coverage proof especially
short.

\section{A ProveIt formalization blueprint}\label{sec:formalization}

The repository snapshot inspected for this report distinguishes formal proofs,
executable certificates, and unformalized research reports \cite{ProveIt}.
The present project can follow that boundary explicitly.

\subsection{Stage I: finite posets}
A small generic library should define finite antichains, chain covers, width,
and budgeted colorings.  The principal theorem can be stated schematically as
\begin{verbatim}
theorem exists_width_budget_coloring_iff
    (P : FinitePoset) (budget : Fin q -> Nat) :
  (exists c, forall i, width (fiber c i) <= budget i) <->
  width P <= sum budget
\end{verbatim}
The upper implication can either import a finite Dilworth theorem from
mathlib or prove the needed chain-partition form through bipartite matching.
The latter produces a reusable algorithmic certificate.

\subsection{Stage II: Boolean lattices}
Formalize the recursive SCD directly.  Its invariants are:
\begin{enumerate}
\item every chain is saturated;
\item bottom rank plus top rank is $n$;
\item chains are pairwise disjoint;
\item their union is the entire Boolean lattice.
\end{enumerate}
The chain count follows by intersecting with the middle layer.  Sperner's
lower bound can be imported if available or proved through the LYM inequality;
for this particular theorem the standard SCD itself also supplies an upper
chain cover while the middle layer supplies the lower antichain.

\subsection{Stage III: set-theoretic local expansion}
In the first-order $\ZF$ development, define the assertion that a function's
fibers contain no $(r+1)$-element antichain.  The restriction to a finite
subcube should then instantiate the finite theorem.  The hull recursion in
\cref{thm:dedekind} requires only finite powersets, least elements of finite
linear orders, and ordinary recursion on $\omega$.

\subsection{Stage IV: antichain transfer}
The complementary-transversal code is elementary and should formalize
independently of cardinal arithmetic.  The dependency graph is:
\[
\boxed{U\times2\hookrightarrow X}
\Longrightarrow
\boxed{\Pow(U)\to\Pow(X)\text{ antichain code}}
\Longrightarrow
\boxed{\Pow(U)\to X\text{ boundedly finite-to-one}}
\]
followed by a finite-to-one map $X\to U$ and Forster's theorem.  This makes
Forster's theorem a named interface rather than a hidden metatheoretic step.

\subsection{Stage V: independent verification}
A Rocq/Coq port of the finite certificate checker would be modest and useful.
The set-theoretic theorem can remain Lean-first until a common encoding of
finite-to-one cardinal comparison is available.  The article should continue
to mark all unformalized infinite dependencies explicitly.

\section{Further research questions and topics}\label{sec:questions}

\begin{enumerate}
\item \textbf{The main $\ZF$ classification.}  Prove or refute
      \cref{conj:BWG} for one fixed $r\geq2$, beginning with $r=2$.

\item \textbf{Permutation-model construction.}  Can a Fraenkel--Mostowski or
      symmetric-extension model contain a Dedekind-infinite, non-duplicable
      $X$ with $\SC_2(X)$?  The map must satisfy the finite expansion bounds
      equivariantly.

\item \textbf{Weak choice strength.}  Determine the weakest standard choice
      principle implying the full finite classification for every $r$.
      Candidates include finite-choice schemes, the Boolean prime ideal
      theorem, and compactness principles sufficient for finite-width chain
      decompositions.

\item \textbf{Canonical chain splitting.}  Does every width-$r$ subfamily of
      a power set admit, in $\ZF$, a decomposition into $r$ chains after
      adjoining a finite amount of structure canonically definable from the
      family?

\item \textbf{Finite-to-one core extraction.}  Does $\SC_r(X)$ itself imply
      the existence of an infinite $U$ with $U\times2\preccurlyeq X$ and
      $X\fto U$?  A positive answer would settle the main conjecture through
      \cref{thm:core}.

\item \textbf{Matching formulation.}  The values
      $f(\{x,c_n\})$ from a definable countable core generate a bipartite graph
      with infinite left degree and bounded right degree.  Identify the exact
      choice principle needed for a two-matching saturating the left side.

\item \textbf{Variable fiber budgets.}  For an infinite target, replace the
      constant $r$ by a function $b:X\to\NN$.  The finite restriction gives
      $W_n\leq\sum_{x\in f[\Pow(S)]}b(x)$.  Which growth or summability
      hypotheses force nonexistence in $\ZF$?

\item \textbf{Stability.}  Classify finite colorings of $\B_n$ using
      $\lceil W_n/r\rceil+o(W_n)$ labels.  Must most middle-layer classes be
      nearly saturated, and must most fibers be close to unions of $r$
      chains?

\item \textbf{Enumeration.}  Count optimal labeled and unlabeled
      $r$-width compressions of $\B_n$.  Even the subclass arising from grouped
      symmetric-chain decompositions has a nontrivial orbit structure.

\item \textbf{Other Peck posets.}  For any finite poset the numerical optimum
      is width divided by $r$, but Boolean lattices offer canonical
      certificates.  Develop equally explicit constructions for products of
      chains, subspace lattices, Young lattices at fixed rank, and divisor
      lattices.

\item \textbf{$q$-ary cubes.}  Replace $\Pow([n])$ by $[q]^n$ under coordinate
      order.  Determine exact self-indexed thresholds when the target has
      size $n$, $qn$, or the size of a prescribed rank.

\item \textbf{Transfinite width bounds.}  Let the allowed fiber width be an
      infinite cardinal $\rho$.  Under which assumptions does an antichain
      code convert a $\rho$-width map into a cardinal inequality strong enough
      to contradict generalized Cantor theorems?

\item \textbf{Reverse mathematics of the finite proof.}  Over weak arithmetic,
      calibrate the strength of the uniform statement that every finite poset
      obeys \cref{thm:budget}, distinguishing existence from an executable
      chain decomposition.

\item \textbf{Formalize Forster.}  Produce a self-contained Lean proof of the
      theorem that a finite-to-one $\Pow(X)\to X$ forces $X$ finite, with an
      explicit audit of every finite-choice step avoided by the argument.

\item \textbf{Definability strength.}  Analyze the sequence produced in
      \cref{thm:dedekind}: can the finite seed be removed, or can one obtain a
      parameter-free countable subset from structural assumptions on $f$?

\item \textbf{Benchmark variants.}  Package the finite theorem, the threshold
      inversion, and the antichain-transfer lemma as separate automatically
      verifiable challenge problems.  A high-end version can ask solvers to
      discover the Dedekind-infinitude recursion without being told the hull
      operator.
\end{enumerate}

\section{Conclusion}\label{sec:conclusion}

Bounded-width power-set compression is a compact problem with three distinct
layers.  The finite layer is exact and universal: width is the only resource,
and a middle antichain plus a grouped chain decomposition gives a complete
certificate.  For Boolean lattices this produces the central-binomial
threshold and its logarithmic inversion.  The classical $r=1$ case is not
merely a finite observation; Glazer's theorem shows that the same cutoff is
absolute in $\ZF$.

For larger $r$, the infinite problem exposes precisely where choice enters.
The local finite theorem still acts inside every finite subcube and forces a
canonical countable subset.  A large complementary-transversal antichain,
combined with Forster's theorem, eliminates every binary-duplicable set and
therefore every well-orderable infinite set.  What remains is not a vague
``pathological set'' possibility but a sharply delimited class of
Dedekind-infinite, non-duplicable objects with unusual finite-multiple
arithmetic.

This combination -- exact finite extremal structure, an objective certificate
surface, a clean $\ZFC$ classification, and a narrow unresolved choiceless
boundary -- makes the problem well suited both to Glazer's research themes
and to ProveIt's formal research program.

\appendix
\section{Exact tables}\label{app:tables}

\begin{table}[H]
\centering
\caption{Central width and least self-indexed defect.}
\label{tab:defects}
\begin{tabular}{r@{\qquad}r@{\qquad}r}
\toprule
$n$ & $W_n=\binom n{\lfloor n/2\rfloor}$ & $\delta(n)=\lceil W_n/n\rceil$\\
\midrule
1 & 1 & 1 \\
2 & 2 & 1 \\
3 & 3 & 1 \\
4 & 6 & 2 \\
5 & 10 & 2 \\
6 & 20 & 4 \\
7 & 35 & 5 \\
8 & 70 & 9 \\
9 & 126 & 14 \\
10 & 252 & 26 \\
11 & 462 & 42 \\
12 & 924 & 77 \\
13 & 1{,}716 & 132 \\
14 & 3{,}432 & 246 \\
15 & 6{,}435 & 429 \\
16 & 12{,}870 & 805 \\
17 & 24{,}310 & 1{,}430 \\
18 & 48{,}620 & 2{,}702 \\
19 & 92{,}378 & 4{,}862 \\
20 & 184{,}756 & 9{,}238 \\
21 & 352{,}716 & 16{,}796 \\
22 & 705{,}432 & 32{,}066 \\
23 & 1{,}352{,}078 & 58{,}786 \\
24 & 2{,}704{,}156 & 112{,}674 \\
\bottomrule
\end{tabular}
\end{table}

\begin{table}[H]
\centering
\caption{Selected exact thresholds $N(r)$.}
\label{tab:thresholds}
\begin{tabular}{r@{\qquad}r|r@{\qquad}r}
\toprule
$r$ & $N(r)$ & $r$ & $N(r)$\\
\midrule
1 & 3 & 26 & 10 \\
2 & 5 & 32 & 10 \\
3 & 5 & 64 & 11 \\
4 & 6 & 128 & 12 \\
5 & 7 & 256 & 14 \\
8 & 7 & 512 & 15 \\
9 & 8 & 1{,}024 & 16 \\
16 & 9 & 4{,}096 & 18 \\
65{,}536 & 23 & 1{,}048{,}576 & 27 \\
\bottomrule
\end{tabular}
\end{table}

\section{Claim and dependency ledger}\label{app:ledger}

\begin{longtable}{L{0.24\linewidth}L{0.28\linewidth}L{0.38\linewidth}}
\toprule
Claim & Status & Dependency or verification\\
\midrule
\endfirsthead
\toprule
Claim & Status & Dependency or verification\\
\midrule
\endhead
Budgeted finite-poset theorem & Proved here & Dilworth's finite chain-decomposition theorem.\\
Exact Boolean formula & Proved here & Sperner's width theorem plus the budget theorem; SCD gives explicit upper certificate.\\
Finite self-indexed classification & Proved here & Immediate from the exact Boolean formula.\\
Monotonicity and asymptotics & Proved here & Exact binomial ratios and Stirling's formula.\\
Glazer chain-fiber classification & Imported & MathOverflow answer of 2 December 2025; not reproved in full.\\
Exact local expansion & Proved here & Restriction to a finite Boolean subcube.\\
Definable Dedekind-infinitude & Proved here & Iterated finite hull with canonical least-preimage orders.\\
Forster theorem & Imported & JSL 68 (2003), 1251--1253.\\
Antichain-transfer theorem & Proved here & Width-to-multiplicity lemma plus Forster.\\
Binary duplicability obstruction & Proved here & Explicit complementary-transversal antichain code.\\
Full $\ZFC$ classification & Proved here & Binary duplicability of infinite well-orderable sets plus finite theorem.\\
Full $\ZF$ classification for $r\geq2$ & Open in this manuscript & Stated as \cref{conj:BWG}; no solution claimed.\\
Historical novelty & Not established & Targeted source search only; independent literature review required.\\
Formal verification & Not claimed & Python checks finite instances; no Lean/Rocq proof accompanies the article.\\
\bottomrule
\end{longtable}

\section{Reproducing the checks}\label{app:reproduce}

From the extracted archive, run
\begin{verbatim}
python3 verification.py
\end{verbatim}
The script rewrites \repo{verification_report.txt} and exits nonzero if any
assertion fails.  No third-party Python package is required.  To rebuild the
article, run
\begin{verbatim}
latexmk -pdf -interaction=nonstopmode -halt-on-error article.tex
\end{verbatim}
or use the PDF skill's \repo{latex_to_pdf.py} wrapper.

\begin{thebibliography}{99}

\bibitem{deBruijn1951}
N.~G. de Bruijn, C. van Ebbenhorst Tengbergen, and D. Kruyswijk,
\newblock On the set of divisors of a number,
\newblock \emph{Nieuw Archief voor Wiskunde} (2) \textbf{23} (1951), 191--193.

\bibitem{Dilworth1950}
R.~P. Dilworth,
\newblock A decomposition theorem for partially ordered sets,
\newblock \emph{Annals of Mathematics} (2) \textbf{51} (1950), 161--166.

\bibitem{Forster2003}
Thomas Forster,
\newblock Finite-to-one maps,
\newblock \emph{Journal of Symbolic Logic} \textbf{68} (2003), no.~4,
1251--1253, doi:10.2178/jsl/1067620184.

\bibitem{FrontierMath}
Elliot Glazer et al.,
\newblock FrontierMath: A benchmark for evaluating advanced mathematical reasoning in AI,
\newblock arXiv:2411.04872, version 7, 23 December 2025.

\bibitem{GlazerYao2026}
Elliot Glazer and Bokai Yao,
\newblock Reflection principles in ZFU,
\newblock arXiv:2602.21970, 25 February 2026.

\bibitem{GreeneKleitman1976}
Curtis Greene and Daniel J. Kleitman,
\newblock Strong versions of Sperner's theorem,
\newblock \emph{Journal of Combinatorial Theory, Series A} \textbf{20}
(1976), 80--88.

\bibitem{HuShen2025}
Xiao Hu and Guozhen Shen,
\newblock Boundedly finite-to-one functions,
\newblock \emph{Logic Journal of the IGPL} \textbf{33} (2025), no.~3,
doi:10.1093/jigpal/jzae130; arXiv:2407.10183v4.

\bibitem{ProveIt}
Vladimir Reshetnikov,
\newblock \emph{ProveIt: machine-checked mathematics in Lean 4 and Rocq/Coq},
\newblock GitHub repository, snapshot
\texttt{b71545625}, 5 October 2026,
\url{https://github.com/VladimirReshetnikov/ProveIt}.

\bibitem{PrussGlazer2025}
Alexander Pruss, Elliot Glazer, Emil Je\v{r}\'abek, and contributors,
\newblock Without AC can there be a function with linearly ordered fibers from the powerset of an infinite set to the set?,
\newblock MathOverflow question 504570 and answers, 2--4 December 2025,
\url{https://mathoverflow.net/questions/504570/}.

\bibitem{Shen2026}
Guozhen Shen,
\newblock Cantor's theorem may fail for finitary partitions,
\newblock \emph{Journal of Symbolic Logic} \textbf{91} (2026), no.~2,
599--616; arXiv:2309.00235.

\bibitem{Sperner1928}
Emanuel Sperner,
\newblock Ein Satz \"uber Untermengen einer endlichen Menge,
\newblock \emph{Mathematische Zeitschrift} \textbf{27} (1928), 544--548.

\end{thebibliography}

\end{document}
